\subsection{Exercises}

In this section, we will provide some exercises to practice the concepts learned in the previous sections.

\subsubsection{Single-Qubit States and Probabilities}\label{sec:single-qubit-states-and-probabilities}

Let's start with a simple exercise to understand single-qubit states and their associated probabilities.

\highspace
\exercise Given:
\begin{equation*}
    \ket{\psi} = \dfrac{1}{2}\ket{0} + \dfrac{\sqrt{3}}{2}\ket{1}
\end{equation*}
Compute:
\begin{equation*}
    P\left(q = 0\right) \quad \text{and} \quad P\left(q = 1\right)
\end{equation*}

\highspace
\solution First of all, we identify the complex probability amplitudes of the state $\ket{\psi}$ (\autopageref{sec:state-space-of-a-single-qubit-system}):
\begin{equation*}
    \ket{\psi} = \alpha\ket{0} + \beta\ket{1} \quad \Rightarrow \quad \ket{\psi} = \dfrac{1}{2}\ket{0} + \dfrac{\sqrt{3}}{2}\ket{1} \quad \Rightarrow \quad \alpha = \dfrac{1}{2}, \quad \beta = \dfrac{\sqrt{3}}{2}
\end{equation*}
To compute the probabilities, we use the formula for the probability of measuring a particular state (\autopageref{eq:magnitude}):
\begin{equation*}
    P\left(q = 0\right) = \left|\alpha\right|^2 = \left|\dfrac{1}{2}\right|^2 = \dfrac{1}{4} = 0.25 = 25\%
\end{equation*}
\begin{equation*}
    P\left(q = 1\right) = \left|\beta\right|^2 = \left|\dfrac{\sqrt{3}}{2}\right|^2 = \dfrac{3}{4} = 0.75 = 75\%
\end{equation*}
As a sanity check, we can verify that the sum of the probabilities equals 1 (the quantum state is normalized):
\begin{equation*}
    P\left(q = 0\right) + P\left(q = 1\right) = 0.25 + 0.75 = 1
\end{equation*}

\highspace
\exercise Given:
\begin{equation*}
    \ket{\psi} = \dfrac{1}{\sqrt{2}}\ket{0} - \dfrac{1}{\sqrt{2}}\ket{1}
\end{equation*}
Compute:
\begin{equation*}
    P\left(q = 0\right) \quad \text{and} \quad P\left(q = 1\right)
\end{equation*}

\highspace
\solution As before, we identify the complex probability amplitudes of the state $\ket{\psi}$:
\begin{equation*}
    \ket{\psi} = \alpha\ket{0} + \beta\ket{1} \quad \Rightarrow \quad \ket{\psi} = \dfrac{1}{\sqrt{2}}\ket{0} - \dfrac{1}{\sqrt{2}}\ket{1} \quad \Rightarrow \quad \alpha = \dfrac{1}{\sqrt{2}}, \quad \beta = -\dfrac{1}{\sqrt{2}}
\end{equation*}
To compute the probabilities, we use the formula for the probability of measuring a particular state:
\begin{equation*}
    P\left(q = 0\right) = \left|\alpha\right|^2 = \left|\dfrac{1}{\sqrt{2}}\right|^2 = \dfrac{1}{2} = 0.5 = 50\%
\end{equation*}
\begin{equation*}
    P\left(q = 1\right) = \left|\beta\right|^2 = \left|-\dfrac{1}{\sqrt{2}}\right|^2 = \dfrac{1}{2} = 0.5 = 50\%
\end{equation*}
As a sanity check, we can verify that the sum of the probabilities equals 1:
\begin{equation*}
    P\left(q = 0\right) + P\left(q = 1\right) = 0.5 + 0.5 = 1
\end{equation*}
